%------------------------------------------------------------------------------%

\usepackage{calculo_varias_variaveis-1}

\title{Regra da Cadeia}

%------------------------------------------------------------------------------%
\begin{document}

\addtitle

%------------------------------------------------------------------------------%
\section{Regra da Cadeia}

%------------------------------------------------------------------------------%
\begin{frame}{Regra da Cadeia em uma variável}
    \[
        h(x) = f\big(g(x)\big)
    \]

    \pause
    \vspace{\baselineskip}
    \[
        \frac{dh}{dx}(x) = \frac{df}{dg}\big(g(x)\big)\frac{dg}{dx}(x)
    \]

    \pause
    \vspace{\baselineskip}
    \[
        h'(x) = f'\big(g(x)\big) g'(x)
    \]
\end{frame}

%------------------------------------------------------------------------------%
\section{Caso 1}

%------------------------------------------------------------------------------%
\begin{frame}{Regra da Cadeia -- Caso 1}
    Uma variável independente e duas variáveis intermediárias
    \vspace{\baselineskip}

    \pause
    Se \(w=f(x, y)\) é diferenciável
    \pause
    e se \(x=x(t)\) e \(y=y(t)\) são deriváveis

    \pause
    \vspace{.5\baselineskip}
    a composta \(w=f\big(x(t), y(t)\big)\) será uma função derivável de $t$ e
    \pause
    \[
        \frac{dw}{dt}
        \;=\; \emph{\D{f}{x}\big(x(t), y(t)\big)}\; \frac{dx}{dt}(t)
        \;+\; \emph{\D{f}{y}\big(x(t), y(t)\big)}\; \frac{dy}{dt}(t)
    \]
\end{frame}

%------------------------------------------------------------------------------%
\framedgraphic{Regra da Cadeia -- Caso 1}{diagrama-caso-1}{}

%------------------------------------------------------------------------------%
%------------------------------------------------------------------------------%
\begin{question}
    Calcule a derivada de \(w=xy\) com relação a $t$ ao longo do caminho
    \[
        x = \cos(t) \qquad
        y = \sin(t)
    \]
    Qual o valor da derivada em \(t=\frac{\pi}{2}\)
\end{question}

%------------------------------------------------------------------------------%
\begin{resolution}{Solução}
  Pela regra da cadeia
  \[
    \frac{dw}{dt} = \D{w}{x}\frac{dx}{dt} + \D{w}{y}\frac{dy}{dt}
  \]
\end{resolution}

%------------------------------------------------------------------------------%
\begin{resolution}{Solução}
  Calculando as derivadas separadamente
  \[
    \pause
    \D{w}{x}
    \pause
    \;=\; \D{}{x} (xy)
    \pause
    \;=\; y\D{x}{x}
    \pause
    \;=\; y
  \]
  \[
    \pause
    \D{w}{y}
    \pause
    \;=\; \D{}{y}(xy)
    \pause
    \;=\; x\D{y}{y}
    \pause
    \;=\; x
  \]

  \[
    \pause
    \frac{dx}{dt}
    \pause
    \;=\; \frac{d}{dt} \cos(t)
    \pause
    \;=\; -\sin(t)
  \]
  \[
    \pause
    \frac{dy}{dt}
    \pause
    \;=\; \frac{d}{dt} \sin(t)
    \pause
    \;=\; \cos(t)
  \]
\end{resolution}

%------------------------------------------------------------------------------%
\begin{resolution}{Solução}
\begin{align*}
    \onslide<+->{\frac{dw}{dt}}
    \onslide<+->{& = \D{w}{x}\frac{dx}{dt} + \D{w}{y}\frac{dy}{dt} \\}
    \onslide<+->{& = \D{}{x}(xy)\;\frac{d}{dt}\big(\cos(t)\big)
               \;+\; \D{}{y}(xy)\;\frac{d}{dt}\big(\sin(t)\big)    \\}
    \onslide<+->{& = y\big(-\sin(t)\big) + x\big(\cos(t)\big)      \\}
    \onslide<+->{& = \big(\sin(t)\big)\big(-\sin(t)\big) + \big(\cos(t)\big)\big(\cos(t)\big) \\}
    \onslide<+->{& = -\sin^2(t) + \cos^2(t) \\}
    \onslide<+->{& = \cos(2t)}
\end{align*}
\end{resolution}

%------------------------------------------------------------------------------%
\begin{resolution}{Solução}
    \begin{align*}
        \onslide<+->{\frac{dw}{dt}\left(\frac{\pi}{2}\right)  }
        \onslide<+->{& = \cos(2t)\evalat{t=\frac{\pi}{2}}{} \\}
        \onslide<+->{& = \cos\left(2\frac{\pi}{2}\right)    \\}
        \onslide<+->{& = \cos(\pi)                          \\}
        \onslide<+->{& = -1}
    \end{align*}
\end{resolution}

%------------------------------------------------------------------------------%


%------------------------------------------------------------------------------%
\begin{frame}{Com Três Variáveis}
    Se \(w=f(x, y, z)\) é diferenciável

    \pause
     e se \(x=x(t)\), \(y=y(t)\) e \(z=z(t)\) são deriváveis

    \pause
    a composta \(w=f\big(x(t), y(t), z(t)\big)\) será uma função derivável de $t$ e

    \pause
    \begin{align*}
        \frac{dw}{dt}
        & = \emph{\D{f}{x}\big(x(t), y(t), z(t)\big)} \; \frac{dx}{dt}(t) \\[3mm]
        & + \emph{\D{f}{y}\big(x(t), y(t), z(t)\big)} \; \frac{dy}{dt}(t) \\[3mm]
        & + \emph{\D{f}{z}\big(x(t), y(t), z(t)\big)} \; \frac{dz}{dt}(t)
    \end{align*}
\end{frame}

\framedgraphic{Regra da Cadeia -- Caso 1}{diagrama-caso-1-3D}{}

%------------------------------------------------------------------------------%
\section{Caso 2}

%------------------------------------------------------------------------------%
\begin{frame}{Regra da Cadeia -- Caso 2 }
    Funções definidas em superfícies
    \vspace{0.5\baselineskip}

    \pause
    Se \(w=f(x, y, z)\) é diferenciável

    \pause
    e se \(x=x(r, s)\), \(y=y(r, s)\) e \(z=z(r, s)\) são deriváveis

    \pause
    a composta \(w=f\big(x(r, s), y(r, s), z(r, s)\big)\)
    será uma função das deriváveis $r$ e $s$

    \pause
    e terá duas derivadas parciais
    \[
        \D{w}{r}
        = \emph{\D{f}{x}} \D{x}{r}
        + \emph{\D{f}{y}} \D{y}{r}
        + \emph{\D{f}{z}} \D{z}{r}
    \]
    \[
        \pause
        \D{w}{s}
        = \emph{\D{f}{x}} \D{x}{s}
        + \emph{\D{f}{y}} \D{y}{s}
        + \emph{\D{f}{z}} \D{z}{s}
    \]
\end{frame}

%------------------------------------------------------------------------------%
\framedgraphic{Regra da Cadeia -- Caso 2}{diagrama-caso-2}{}

%------------------------------------------------------------------------------%
%------------------------------------------------------------------------------%
\begin{question}
    Expresse \(\D{w}{r}\) e \(\D{w}{s}\) em termos de $r$ e $s$
    \[
        w = x + 2y + z^2    \qquad
        x = \frac{r}{s}     \qquad
        y = r^2 + \ln(s)    \qquad
        z = 2r
    \]
\end{question}

%------------------------------------------------------------------------------%
\begin{resolution}{Solução}
  Pela regra da cadeia

  \begin{align*}
    \D{w}{r} & = \D{w}{x}\D{x}{r} + \D{w}{y}\D{y}{r} + \D{w}{z}\D{z}{r}  \\[3mm]
    \D{w}{s} & = \D{w}{x}\D{x}{s} + \D{w}{y}\D{y}{s} + \D{w}{z}\D{z}{s}
  \end{align*}
\end{resolution}

%------------------------------------------------------------------------------%
\begin{resolution}{Solução}
  Calculando as derivadas separadamente

  \[
    \pause
    \D{w}{x}
    \pause
    \;=\; \D{}{x}\left(x + 2y + z^2\right)
    \pause
    \;=\; 1
  \]

  \[
    \pause
    \D{w}{y}
    \pause
    \;=\; \D{}{y}\left(x + 2y + z^2\right)
    \pause
    \;=\; 2
  \]

  \[
    \pause
    \D{w}{z}
    \pause
    \;=\; \D{}{z}\left(x + 2y + z^2\right)
    \pause
    \;=\; 2z
  \]
\end{resolution}

%------------------------------------------------------------------------------%
\begin{resolution}{Solução}
\begin{columns}
\column{0.4\linewidth}
  \[
    \D{x}{r}
    \pause
    \;=\; \D{}{r}\left(\frac{r}{s}\right)
    \pause
    \;=\; \frac{1}{s}
  \]

  \[
    \pause
    \D{y}{r}
    \pause
    \;=\; \D{}{r}\left(r^2 + \ln(s)\right)
    \pause
    \;=\; 2r
  \]

  \[
    \pause
    \D{z}{r}
    \pause
    \;=\; \D{}{r}\left(2r\right)
    \pause
    \;=\; 2
  \]
\column{0.6\linewidth}
  \[
    \pause
    \D{x}{s}
    \pause
    \;=\; \D{}{s}\left(\frac{r}{s}\right)
    \pause
    \;=\; \D{}{s}\left(rs^{-1}\right)
    \pause
    \;=\; -\frac{r}{s^2}
  \]

  \[
    \pause
    \D{y}{s}
    \pause
    \;=\; \D{}{s}\left(r^2 + \ln(s)\right)
    \pause
    \;=\; \frac{1}{s}
  \]

  \[
    \pause
    \D{z}{s}
    \pause
    \;=\; \D{}{s}\left(2r\right)
    \pause
    \;=\; 0
  \]
\end{columns}
\end{resolution}

%------------------------------------------------------------------------------%
\begin{resolution}{Solução}
    \begin{align*}
        \onslide<+->{\D{w}{r}}
        \onslide<+->{& = \D{w}{x}\D{x}{r} + \D{w}{y}\D{y}{r} + \D{w}{z}\D{z}{r}  \\}
        \onslide<+->{& = 1\D{x}{r} + 2\D{y}{r} + 2z\D{z}{r}  \\}
        \onslide<+->{& = 1\frac{1}{s} + 2(2r) + 2z2          \\}
        \onslide<+->{& = \frac{1}{s} + 4r + 4z               \\}
        \onslide<+->{& = \frac{1}{s} + 4r + 4(2r)            \\}
        \onslide<+->{& = \frac{1}{s} + 12r }
    \end{align*}
\end{resolution}

%------------------------------------------------------------------------------%
\begin{resolution}{Solução}
    \begin{align*}
        \onslide<+->{\D{w}{s}}
        \onslide<+->{& = \D{w}{x}\D{x}{s} + \D{w}{y}\D{y}{s} + \D{w}{z}\D{z}{s}   \\}
        \onslide<+->{& = 1\D{x}{s} + 2\D{y}{s} + 2z\D{z}{s}                       \\}
        \onslide<+->{& = 1\left(\frac{-r}{s^2}\right) + 2\frac{1}{s} + 2z\times 0 \\}
        \onslide<+->{& = -\frac{r}{s^2} + \frac{2}{s}}
    \end{align*}
\end{resolution}

%------------------------------------------------------------------------------%


%------------------------------------------------------------------------------%
\section{Exemplos}

%------------------------------------------------------------------------------%
\begin{question}
  Expresse \(\D{w}{r}\) e \(\D{w}{s}\) em termos de $r$ e $s$
  \[
    w = x^2 + y^2  \qquad
    x = r^3 - s^2  \qquad
    y = r^2 + s^3
  \]
\end{question}

%------------------------------------------------------------------------------%
\begin{resolution}{Solução}
  Pela regra da cadeia
  \begin{align*}
    \D{w}{r} & = \D{w}{x}\D{x}{r} + \D{w}{y}\D{y}{r}  \\
    \D{w}{s} & = \D{w}{x}\D{x}{s} + \D{w}{y}\D{y}{s}
  \end{align*}
\end{resolution}

%------------------------------------------------------------------------------%
\begin{resolution}{Solução}
\begin{columns}
\column{0.5\linewidth}
Calculando as derivadas

  \[
    \pause
    \D{w}{x}
    \pause
    \;=\; \D{}{x}\left(x^2 + y^2\right)
    \pause
    \;=\; 2x
  \]

  \[
    \pause
    \D{w}{y}
    \pause
    \;=\; \D{}{y}\left(x^2 + y^2\right)
    \pause
    \;=\; 2y
  \]
\column{0.5\linewidth}
  \[
    \pause
    \D{x}{r}
    \pause
    \;=\; \D{}{r}\left(r^3 - s^2\right)
    \pause
    \;=\; 3r^2
  \]

  \[
    \pause
    \D{y}{r}
    \pause
    \;=\; \D{}{r}\left(r^2 + s^3\right)
    \pause
    \;=\; 2r
  \]

  ~
  \[
    \pause
    \D{x}{s}
    \pause
    \;=\; \D{}{s}\left(r^3 - s^2\right)
    \pause
    \;=\; -2s
  \]

  \[
    \pause
    \D{y}{s}
    \pause
    \;=\; \D{}{s}\left(r^2 + s^3\right)
    \pause
    \;=\; 3s^2
  \]
\end{columns}
\end{resolution}

%------------------------------------------------------------------------------%
\begin{resolution}{Solução}
  \begin{align*}
    \onslide<+->{\D{w}{r}}
    \onslide<+->{& = \D{w}{x}\D{x}{r} + \D{w}{y}\D{y}{r}  \\}
    \onslide<+->{& = 2x\D{x}{r} + 2y\D{y}{r}              \\}
    \onslide<+->{& = 2x3r^2 + 2y2r                        \\}
    \onslide<+->{& = 6(r^3 - s^2)r^2 + 4(r^2 + s^3)r      \\}
    \onslide<+->{& = 6r^5 - 6s^2r^2 +4r^3 + 4s^3r         \\}
    \onslide<+->{& = 6r^5 + 4r^3 - 6r^2s^2 + 4rs^3}
  \end{align*}
\end{resolution}

%------------------------------------------------------------------------------%
\begin{resolution}{Solução}
  \begin{align*}
    \onslide<+->{\D{w}{s}}
    \onslide<+->{& = \D{w}{x}\D{x}{s} + \D{w}{y}\D{y}{s}  \\}
    \onslide<+->{& = 2x\D{x}{s} + 2y\D{y}{s}              \\}
    \onslide<+->{& = 2x\left(-2s\right) + 2y3s^2          \\}
    \onslide<+->{& = -4(r^3 - s^2)s + 6(r^2 + s^3)s^2     \\}
    \onslide<+->{& = -4r^3s +4s^3 +6r^2s^2 + 6s^5}
  \end{align*}
\end{resolution}

%------------------------------------------------------------------------------%

%------------------------------------------------------------------------------%
\begin{question}
  Considere \(z = \sin\left(x^2 + y^2\right)\)
  onde \(x = st\) e \(y = 2s + t\)

  Encontre \(\frac{\partial z}{\partial s}\)
\end{question}

%------------------------------------------------------------------------------%
\begin{resolution}{Solução}
  Pela regra da cadeia
  \[
    \D{z}{s} = \D{z}{x}\D{x}{s} + \D{z}{y}\D{y}{s}
  \]
\end{resolution}

%------------------------------------------------------------------------------%
\begin{resolution}{Solução}
  \[
    \D{z}{x}
    \pause
    \;=\; \D{}{x} \sin\left(x^2 + y^2\right)
    \pause
    \;=\; \cos\left(x^2 + y^2\right) \D{}{x}\left(x^2 + y^2\right)
    \pause
    \;=\; \cos(x^2 + y^2) 2x
  \]
  \[
    \pause
    \D{z}{y}
    \pause
    \;=\; \D{}{y} \sin\left(x^2 + y^2\right)
    \pause
    \;=\; \cos\left(x^2 + y^2\right) \D{}{y}\left(x^2 + y^2\right)
    \pause
    \;=\; \cos(x^2 + y^2) 2y
  \]

  \[
    \pause
    \D{x}{s}
    \pause
    \;=\; \D{}{s} \left(st\right)
    \pause
    \;=\; t
  \]
  \[
    \pause
    \D{y}{s}
    \pause
    \;=\; \D{}{s} \left(2s+t\right)
    \pause
    \;=\; 2
  \]
\end{resolution}

%------------------------------------------------------------------------------%
\begin{resolution}{Solução}
\begin{align*}
  \onslide<+->{\D{z}{s}}
  \onslide<+->{& = \D{z}{x}\D{x}{s} + \D{z}{y}\D{y}{s}                              \\}
  \onslide<+->{& = 2x\cos\left(x^2 + y^2\right)\D{x}{s} + 2y\cos\left(x^2 + y^2\right)\D{y}{s} \\}
  \onslide<+->{& = 2x\cos\left(x^2 + y^2\right) t
    + 2y\cos\left(x^2 + y^2\right) 2                                                \\}
  \onslide<+->{& = \left(2xt + 4y \right) \cos\left(x^2 + y^2\right)                \\}
  \onslide<+->{& = \left[ 2(st)t + 4(s+t) \right] \cos\left((st)^2 + (s+t)^2\right) \\}
  \onslide<+->{& = \left( 2st^2 + 4s +4t \right) \cos\left(s^2t^2 + (s+t)^2\right)}
\end{align*}
\end{resolution}

%------------------------------------------------------------------------------%

%------------------------------------------------------------------------------%
\begin{question}
Seja \(w = e^{xyz}\) onde \(x = r^2 + 4s\), \(y = 2r - 3s\), \(z = rs\)

Encontre \(\D{w}{r}\)
\end{question}

%------------------------------------------------------------------------------%
\begin{resolution}{Solução}
  Pela regra da cadeia
  \[
      \D{w}{r}
    = \D{w}{x}\D{x}{r}
    + \D{w}{y}\D{y}{r}
    + \D{w}{z}\D{z}{r}
  \]
\end{resolution}

%------------------------------------------------------------------------------%
\begin{resolution}{Solução}
\begin{columns}
\column{0.5\linewidth}
  \[
    \D{w}{x}
    \pause
    \;=\; \D{}{x} e^{xyz}
    \pause
    \;=\; yz e^{xyz}
  \]
  \[
    \pause
    \D{w}{y}
    \pause
    \;=\; \D{}{y} e^{xyz}
    \pause
    \;=\; xz e^{xyz}
  \]
  \[
    \pause
    \D{w}{z}
    \pause
    \;=\; \D{}{z} e^{xyz}
    \pause
    \;=\; xy e^{xyz}
  \]
\column{0.5\linewidth}
  \[
    \pause
    \D{x}{r}
    \pause
    \;=\; \D{}{r}\left(r^2 + 4s\right)
    \pause
    \;=\; 2r
  \]
  \[
    \pause
    \D{y}{r}
    \pause
    \;=\; \D{}{r}\left(2r - 3s\right)
    \pause
    \;=\; 2
  \]
  \[
    \pause
    \D{z}{r}
    \pause
    \;=\; \D{}{r}\left(rs\right)
    \pause
    \;=\; s
  \]
\end{columns}
\end{resolution}

%------------------------------------------------------------------------------%
\begin{resolution}{Solução}
\begin{align*}
  \onslide<+->{\D{w}{r}}
  \onslide<+->{& = \D{w}{x}\D{x}{r} + \D{w}{y}\D{y}{r} + \D{w}{z}\D{z}{r}       \\}
  \onslide<+->{& = yz e^{xyz}\D{x}{r} + xz e^{xyz}\D{y}{r} + xy e^{xyz}\D{z}{r} \\}
  \onslide<+->{& = e^{xyz}\left(yz \D{x}{r} + xz \D{y}{r} + xy \D{z}{r}\right)  \\}
  \onslide<+->{& = e^{xyz}\left(yz2r + xz2 + xys \right)                        \\}
  \onslide<+->{& = e^{xyz}\left[2(2r - 3s)(rs)r + 2(r^2 + 4s)(rs) + (r^2 + 4s)(2r - 3s)s \right] \\}
  \onslide<+->{& = e^{xyz}\left[2r^2s(2r - 3s) + 2rs(r^2 + 4s) + s(r^2 + 4s)(2r - 3s) \right]}
\end{align*}
\end{resolution}

%------------------------------------------------------------------------------%

%------------------------------------------------------------------------------%
\begin{question}
Dada
\(
  u = \ln\left(x^2 + y^2 + z^2\right)
\)
com \(x = e^t\), \(y = e^{-t}\), \(z = t^2\)

Encontre \(\frac{du}{dt}\)
\end{question}

%------------------------------------------------------------------------------%
\begin{resolution}{Solução}
  Pela regra da cadeia
  \[
    \frac{du}{dt}
    = \D{u}{x}\frac{dx}{dt}
    + \D{u}{y}\frac{dy}{dt}
    + \D{u}{z}\frac{dz}{dt}
  \]
\end{resolution}

%------------------------------------------------------------------------------%
\begin{resolution}{Solução}
  \[
    \D{u}{x}
    \pause
    \;=\; \ln'\left(x^2 + y^2 + z^2\right) \D{}{x}\left(x^2 + y^2 + z^2\right)
    \pause
    \;=\; \frac{2x}{x^2 + y^2 + z^2}
  \]

  \[
    \pause
    \D{u}{y}
    \pause
    \;=\; \ln'\left(x^2 + y^2 + z^2\right) \D{}{y}\left(x^2 + y^2 + z^2\right)
    \pause
    \;=\; \frac{2y}{x^2 + y^2 + z^2}
  \]

  \[
    \pause
    \D{u}{z}
    \pause
    \;=\; \ln'\left(x^2 + y^2 + z^2\right) \D{}{z}\left(x^2 + y^2 + z^2\right)
    \pause
    \;=\; \frac{2z}{x^2 + y^2 + z^2}
  \]
\end{resolution}

%------------------------------------------------------------------------------%
\begin{resolution}{Solução}
  \[
    \frac{dx}{dt}
    \pause
    \;=\; \frac{de^t}{dt}
    \pause
    \;=\; e^t
  \]

  \[
    \pause
    \frac{dy}{dt}
    \pause
    \;=\; \frac{de^{-t}}{dt}
    \pause
    \;=\; -e^{-t}
  \]

  \[
    \pause
    \frac{dz}{dt}
    \pause
    \;=\; \frac{dt^2}{dt}
    \pause
    \;=\; 2t
  \]
\end{resolution}

%------------------------------------------------------------------------------%
\begin{resolution}{Solução}
  \begin{align*}
    \onslide<+->{\frac{du}{dt}}
    \onslide<+->{& = \D{u}{x}\frac{dx}{dt} + \D{u}{y}\frac{dy}{dt} + \D{u}{z}\frac{dz}{dt}  \\}
    \onslide<+->{& = \frac{2x}{x^2 + y^2 + z^2}\frac{dx}{dt}
      + \frac{2y}{x^2 + y^2 + z^2}\frac{dy}{dt}
      + \frac{2z}{x^2 + y^2 + z^2}\frac{dz}{dt}  \\}
    \onslide<+->{& = \frac{2}{x^2 + y^2 + z^2}\left(
          x\frac{dx}{dt}
        + y\frac{dy}{dt}
        + z\frac{dz}{dt}
        \right) \\}
    \onslide<+->{& = \frac{2}{x^2 + y^2 + z^2}\left(xe^t -ye^{-t} + 2zt \right) \\}
    \onslide<+->{& = \frac{2e^{2t} -2e^{-2t} + 4t^3}{x^2 + y^2 + z^2}}
  \end{align*}
\end{resolution}

%------------------------------------------------------------------------------%


%------------------------------------------------------------------------------%
\section{Lista Mínima}

%------------------------------------------------------------------------------%
\begin{frame}{Lista Mínima}
  Cálculo Vol. 2 do Thomas 12\textsuperscript{a} ed. -- Seção 14.4
  \begin{enumerate}
    \item Estudar o texto da seção
    \item Resolver os exercícios: 3, 5, 6, 10, 12, 14, 17
  \end{enumerate}

  \vfill
  Atenção: A prova é baseada no livro, não nas apresentações
\end{frame}

%------------------------------------------------------------------------------%
\end{document}
%------------------------------------------------------------------------------%
